% SPDX-License-Identifier: CC-BY-SA-4.0
% Author: Matthieu Perrin

\begingroup

\SetKwFunction{csend}{causal-send}
\SetKwFunction{crec}{causal-receive}

\begin{exercice}[Communication causale point-à-point]

  Le but de cet exercice est de concevoir un service de communication
  \emph{causale point-à-point}, permettant aux processus de s'échanger des
  messages à l'aide des primitives \csend et \crec.

  \begin{questions}
  \item Considérons l'exécution suivante.

    \begin{center}
      \begin{tikzpicture}[y=7mm]
        \draw[process] (0, 3) node[left]{$p_1$} to (7, 3);
        \draw[process] (0, 2) node[left]{$p_2$} to (7, 2);
        \draw[process] (0, 1) node[left]{$p_3$} to (7, 1);

        \draw[message] (2, 3) to[out=-60, in=180] (4, 1.5)
        to[in=150, out=0] (5.5, 1) node[above right] {$m_1$};
        \path[message] (3, 3) edge (3.5, 2) node[above] {$m_2$};
        \path[message] (4, 2) edge (4.5, 1) node[above] {$m_3$};
      \end{tikzpicture}
    \end{center}

    Expliquer pourquoi la réception de $m_3$ avant celle de $m_1$ par $p_3$
    constitue une violation de la causalité.
    En déduire une spécification de la propriété d'ordre que doit garantir
    le service.

  \item Supposons que chaque processus $p_i$ maintienne une horloge vectorielle
    $V_i[1..n]$, transmise avec chacun de ses messages.
    Construire deux exécutions dans lesquelles un processus reçoit un message
    portant la même horloge vectorielle, mais où les messages qui doivent avoir
    été reçus auparavant par ce processus ne sont pas les mêmes.

  \item Proposer une structure de données permettant à chaque processus de
    représenter les dépendances causales entre les communications point-à-point.
    Préciser la signification de chacune de ses composantes.

  \item Proposer un algorithme implémentant les primitives \csend et \crec
    à l'aide de cette structure.
    Préciser les informations transmises avec chaque message ainsi que la
    condition permettant de délivrer un message reçu.

  \item Vérifier sur l'exécution précédente que $m_3$ ne peut pas être délivré
    par $p_3$ avant $m_1$.

  \end{questions}

\end{exercice}

\endgroup
\endinput
